\section{Why the cubic scale is natural}
\label{sec:exponent-accounting}
%======================================================================

Before asking for a proof, we check that every scale is dimensionally consistent.  This section explains why $N^3$, rather than the coherent $N^5$ ceiling, is the natural target.

The block increment
\begin{equation}
 d_n=S_n-S_{n-1}
 \label{eq:d-n}
\end{equation}
satisfies the trivial bound $|d_n|\le2n+1$.  Expanding the cumulative energy gives
\begin{equation}
 \sum_{n\le N}|S_n|^2
 =
 \sum_{j,k\le N}
 \bigl(N-\max\{j,k\}+1\bigr)d_jd_k.
 \label{eq:triangular-energy}
\end{equation}
The fully coherent ceiling is
\begin{equation}
 N^5
 =
 \underbrace{N}_{\text{prefix multiplicity}}
 \times
 \underbrace{N^2}_{\text{block pairs}}
 \times
 \underbrace{N^2}_{\text{raw pair amplitude}}.
 \label{eq:N5-accounting}
\end{equation}

The squared map provides the exact missing scale:
\begin{equation}
 J(c,q)=c^2+q^2\ge2cq\asymp n^2.
 \label{eq:two-power-J}
\end{equation}
Thus a unit area of factor combinations is diluted by $n^{-2}$ in the squared plane.  Equivalently, each coordinate direction is stretched by order $n$.  If the localized kernel respects this conformal dilution, then the effective pair amplitude becomes order one:
\begin{equation}
 N^3
 =
 \underbrace{N}_{\text{prefix multiplicity}}
 \times
 \underbrace{N^2}_{\text{block pairs}}
 \times
 \underbrace{1}_{\text{squared-space interaction}}.
 \label{eq:N3-accounting}
\end{equation}

The low-height theorem proves this reduction directly in the region $|Y|\lesssim n$.  The pointwise high-height conjecture asks for the corresponding signed cancellation in the remaining region.
